\clearpage\SourceLeaf{124}%
\pdfbookmark[1]{Der Zodiacus der Juno}{zodiacus1805}
\begin{center}
{\small DER ZODIACUS DER JUNO.\par}
\medskip
{\small (Ein Nachtrag zu dem Aufsatze im August-Hefte der \textit{M. C.} 1804, [S.\,106 d.\,B.])\par}
\medskip\hrule\medskip
{\small Monatliche Correspondenz zur Beförderung der Erd- und Himmels-Kunde,\par
herausgegeben vom Freiherrn \textsc{von Zach}. März 1805.\par}
\medskip\hrule
\end{center}\medskip
Obgleich die Juno nur erst eine kurze Zeit hindurch beobachtet worden ist,
so scheinen doch die Elemente ihrer Bahn bereits einen hinlänglichen Grad
von Genauigkeit erlangt zu haben, um zum Behuf des von dem verdienstvollen
Entdecker dieses Planeten zu hoffenden Atlasses die Grenzen der Zone, worin er
uns erscheinen kann, abzustecken. Ich habe daher diese Arbeit, auf Ersuchen
meines Freundes \textsc{Harding}, um so lieber übernommen, da gerade die kleine
lichtschwache Juno in weniger günstigen Lagen, als sie dieses Jahr hatte, von
allen drei neuen Planeten am schwersten zu beobachten, und also detaillirter
Sternkarten am meisten bedürftig sein wird.

Meine IV. Elemente der Juno scheinen nach meinen letzten Beobachtungen
noch so gut mit dem Laufe derselben übereinzustimmen, dass ich noch keine
zuverlässige neue Verbesserung der Bahn zu unternehmen im Stande sein würde;
ich habe sie daher bei meinen Rechnungen zum Grunde gelegt, und in den
Zeichen des erwähnten Aufsatzes angenommen:
\[
\begin{aligned}
e&=0.25684\\
k&=2.49630\\
g&=62^0\;19'\;35''\\
e'&=0.01679\\
k'&=0.999718\\
g'&=108^0\;30'\;1''
\end{aligned}
\]
\SourceLeaf{125}%
Hieraus fand ich folgende Bedingungs-Gleichung:
\[
\cos t'=0.400477\cos t-0.053103
\]

Die Werthe von $\cos t'$ liegen also zwischen $+0.347374$ und $-0.453580$,
folglich die von $t'$ einerseits zwischen $69^0\;40'\;23''$ und $116^0\;58'\;25''$; und
andererseits zwischen $234^0\;5'\;47''$ und $281^0\;23'\;49''$, also die heliocentrischen Örter
der Erde, wo Juno in den Limiten erscheinen kann, von $240^0\;45'$ bis $288^0\;3'$
und von $54^0\;6'$ bis $101^0\;24'$ Länge; die erstern fallen etwa vom 22.\,Mai bis 10.\,Julius,
wo Juno einmal in der nördlichen Limite erscheinen muss, die andern
vom 16.\,November bis 2.\,Januar, wo sie einmal an die südliche Grenze kommt.

Hienach wurden nun, gerade so wie bei der Ceres und Pallas, $36$ Punkte
in der nördlichen, und eben so viele in der südlichen Grenze bestimmt, wobei
zur unmittelbaren Berechnung der Rectascensionen und Declinationen folgende
Formeln und Constanten gebraucht werden:
\[
\begin{array}{rcrrr}
a&=&87^0&59'&22''\\
b&=&100&32&38\\
c&=&10&44&17\\
A&=&261&18&7\\
B&=&171&40&35\\
C&=&160&38&2
\end{array}
\]
\[
v=\text{wahre Anomalie der Juno}
\]
\[
\begin{aligned}
x&=\frac{\alpha\sin(v+323^0\;37'\;41'')}{1-e\cos v}\\[8pt]
y&=\frac{\ObservedBeta\sin(v+234^0\;0'\;9'')}{1-e\cos v}\\[8pt]
z&=\frac{\gamma\sin(v+222^0\;57'\;37'')}{1-e\cos v}
\end{aligned}
\]
\[
\begin{aligned}
\alpha&=k\sin a; &\log\alpha&=0.397030\\
\ObservedBeta&=k\sin b; &\log\ObservedBeta&=0.389902\\
\gamma&=k\sin c; &\log\gamma&=9.667556
\end{aligned}
\]

Zwischen diese $72$ Punkte wurden $144$ andere eingeschaltet, wo die Rectascensionen
von $5$ zu $5$ Graden zunehmen, und die in folgender Tafel dargestellt werden:
\SourceLeaf{126}
\begin{center}
\begingroup\small\setlength{\tabcolsep}{3pt}
\begin{tabular}{r|rrl|rrl||r|rrl|rrl}
\multicolumn{14}{c}{\textit{Zodiacus} der Juno.}\\[4pt]
&\multicolumn{6}{c||}{Declination}&&\multicolumn{6}{c}{Declination}\\
AR.&\multicolumn{3}{c|}{\shortstack{der nördlichen\\Grenze}}&\multicolumn{3}{c||}{\shortstack{der südlichen\\Grenze}}&AR.&\multicolumn{3}{c|}{\shortstack{der nördlichen\\Grenze}}&\multicolumn{3}{c}{\shortstack{der südlichen\\Grenze}}\\\hline
$0^0$&$4^0$&$23'$&nördl.&$10^0$&$38'$&südl.&$180^0$&$6^0$&$59'$&nördl.&$3^0$&$9'$&südl.\\
$5$&$5$&$13$&&$10$&$8$&&$185$&$6$&$5$&&$3$&$47$&\\
$10$&$6$&$3$&&$9$&$36$&&$190$&$5$&$11$&&$4$&$27$&\\
$15$&$6$&$51$&&$9$&$3$&&$195$&$4$&$17$&&$5$&$8$&\\
$20$&$7$&$39$&&$8$&$29$&&$200$&$3$&$23$&&$5$&$50$&\\
$25$&$8$&$26$&&$7$&$54$&&$205$&$2$&$30$&&$6$&$32$&\\
$30$&$9$&$11$&&$7$&$19$&&$210$&$1$&$38$&&$7$&$15$&\\
$35$&$9$&$54$&&$6$&$43$&&$215$&$0$&$49$&&$7$&$58$&\\
$40$&$10$&$34$&&$6$&$7$&&$220$&$0$&$11$&&$8$&$40$&\\
$45$&$11$&$13$&&$5$&$32$&&$225$&$0$&$44$&südl.&$9$&$21$&\\
$50$&$11$&$49$&&$4$&$56$&&$230$&$1$&$26$&&$10$&$0$&\\
$55$&$12$&$23$&&$4$&$21$&&$235$&$2$&$3$&&$10$&$38$&\\
$60$&$12$&$53$&&$3$&$47$&&$240$&$2$&$38$&&$11$&$13$&\\
$65$&$13$&$21$&&$3$&$14$&&$245$&$3$&$8$&&$11$&$45$&\\
$70$&$13$&$45$&&$2$&$43$&&$250$&$3$&$34$&&$12$&$15$&\\
$75$&$14$&$7$&&$2$&$13$&&$255$&$3$&$55$&&$12$&$42$&\\
$80$&$14$&$25$&&$1$&$45$&&$260$&$4$&$12$&&$13$&$5$&\\
$85$&$14$&$40$&&$1$&$19$&&$265$&$4$&$24$&&$13$&$26$&\\
$90$&$14$&$51$&&$0$&$56$&&$270$&$4$&$32$&&$13$&$43$&\\
$95$&$14$&$58$&&$0$&$35$&&$275$&$4$&$34$&&$13$&$56$&\\
$100$&$15$&$1$&&$0$&$17$&&$280$&$4$&$32$&&$14$&$7$&\\
$105$&$15$&$1$&&$0$&$2$&&$285$&$4$&$25$&&$14$&$14$&\\
$110$&$14$&$55$&&$0$&$10$&nördl.&$290$&$4$&$14$&&$14$&$17$&\\
$115$&$14$&$46$&&$0$&$19$&&$295$&$3$&$58$&&$14$&$18$&\\
$120$&$14$&$32$&&$0$&$24$&&$300$&$3$&$38$&&$14$&$16$&\\
$125$&$14$&$14$&&$0$&$26$&&$305$&$3$&$14$&&$14$&$11$&\\
$130$&$13$&$52$&&$0$&$24$&&$310$&$2$&$45$&&$14$&$2$&\\
$135$&$13$&$26$&&$0$&$19$&&$315$&$2$&$14$&&$13$&$52$&\\
$140$&$12$&$56$&&$0$&$9$&&$320$&$1$&$38$&&$13$&$39$&\\
$145$&$12$&$22$&&$0$&$4$&südl.&$325$&$1$&$0$&&$13$&$24$&\\
$150$&$11$&$44$&&$0$&$20$&&$330$&$0$&$19$&&$13$&$6$&\\
$155$&$11$&$3$&&$0$&$40$&&$335$&$0$&$25$&nördl.&$12$&$46$&\\
$160$&$10$&$19$&&$1$&$4$&&$340$&$1$&$9$&&$12$&$24$&\\
$165$&$9$&$32$&&$1$&$31$&&$345$&$1$&$57$&&$12$&$0$&\\
$170$&$8$&$43$&&$2$&$1$&&$350$&$2$&$45$&&$11$&$35$&\\
$175$&$7$&$52$&&$2$&$34$&&$355$&$3$&$34$&&$11$&$7$&\\
$180$&$6$&$59$&&$3$&$9$&&$360$&$4$&$23$&&$10$&$38$&\\
\end{tabular}
\endgroup
\end{center}
